% Lecture 3 notes: \input by notes.tex, \subimport'ed by the semester book.
% Exam questions (экзамен): 5–7 (теорема о коллизиях, коллизионный алгоритм для DLP, орбиты,
% ρ-алгоритм Полларда с циклом Флойда), 8 (ρ-алгоритм с выделенными точками).
\begin{lecture}{3}{Алгоритмы поиска коллизий. \texorpdfstring{$\rho$}{ρ}-алгоритм Полларда}

\section{Зачем искать коллизии}

Пусть $G = \langle g \rangle$ "--- циклическая группа порядка $N$ (например, подгруппа
в $\Znx{p}$ или группа точек эллиптической кривой) и $h \in G$. Задача дискретного
логарифмирования (DLP) "--- найти $x$ с $g^x = h$. Стойкость протокола Диффи"--~Хеллмана
и схемы Эль-Гамаля (лекция~1) опирается на трудность этой задачи.

Алгоритм Шенкса (лекция~1) решает DLP за $\bigO(\sqrt N)$ групповых операций, но требует
хранить таблицу из $\approx \sqrt N$ элементов группы. При $N \approx 2^{80}$ это $2^{40}$
элементов "--- терабайты памяти, и именно память, а не время, становится узким местом.
Алгоритм Полига"--~Хеллмана (лекция~2) сводит задачу к подгруппам простого порядка $q \mid N$,
поэтому в криптографии $N$ выбирают так, чтобы у него был большой простой делитель, и
главный вопрос "--- как решать DLP в группе \emph{простого} порядка $q$.

В этой лекции мы увидим, что
\begin{itemize}
  \item \emph{время} $\bigO(\sqrt N)$ "--- это цена поиска \emph{коллизии}, т.\,е. совпадения
    двух случайных элементов множества из $N$ элементов («парадокс дней рождения»);
  \item коллизию можно искать почти без памяти: $\rho$-алгоритм Полларда~\cite{pollard-rho-dlp}
    решает DLP за $\bigO(\sqrt N)$ операций и $\bigO(1)$ памяти;
  \item модификация с выделенными точками~\cite{van-oorschot-wiener} позволяет распараллелить
    поиск на $m$ вычислителей с линейным ускорением "--- так решают рекордные задачи ECDLP.
\end{itemize}
Все эти алгоритмы \emph{универсальны} (generic): они используют только умножение в группе
и сравнение элементов, поэтому одинаково работают в $\Znx{p}$ и на эллиптических кривых
(лекция~9). Для универсальных алгоритмов оценка $\sqrt N$ неулучшаема (теорема Шупа, см.
\cite[гл.~13]{galbraith}).

\section{Парадокс дней рождения и теорема о коллизиях}

Всюду в этом разделе элементы выбираются из множества $S$, $\abs{S} = N$, \emph{независимо
и равновероятно} (с возвращением). \emph{Коллизия} "--- совпадение двух выбранных элементов.

\begin{lemma}\label{l3:exp}
$1 - t \le e^{-t}$ при всех $t \in \RR$.
\end{lemma}

\begin{proof}
Функция $\varphi(t) = e^{-t} - 1 + t$ выпукла ($\varphi'' = e^{-t} > 0$) и
$\varphi(0) = \varphi'(0) = 0$, поэтому $t = 0$ "--- точка минимума и $\varphi \ge 0$.
\end{proof}

\begin{theorem}[один список]\label{l3:birthday}
Пусть $a_1, \dots, a_n \in S$ выбраны независимо и равновероятно. Тогда
\[
  \Pr(\exists\, i < j:\ a_i = a_j)
  = 1 - \prod_{i=1}^{n-1}\Bigl(1 - \frac iN\Bigr) \ \ge\ 1 - e^{-n(n-1)/2N}.
\]
\end{theorem}

\begin{proof}
Все элементы различны, если $a_2 \ne a_1$, $a_3 \notin \{a_1, a_2\}$ и т.\,д. При условии,
что $a_1, \dots, a_i$ различны, элемент $a_{i+1}$ не совпадает ни с одним из них
с вероятностью $1 - i/N$ (он не зависит от предыдущих). Перемножая условные вероятности,
получаем $\Pr(\text{все различны}) = \prod_{i=1}^{n-1} (1 - i/N)$. По
лемме~\ref{l3:exp} это не больше $\prod_{i=1}^{n-1} e^{-i/N} = e^{-\sum i / N} =
e^{-n(n-1)/2N}$.
\end{proof}

При $n \ll N$ неравенство почти обращается в равенство: $\ln(1 - t) = -t + \bigO(t^2)$,
и ошибка в показателе $\sum_{i<n} \bigO(i^2/N^2) = \bigO(n^3/N^2)$ мала при $n = \bigO(\sqrt N)$.

\begin{example}[парадокс дней рождения]
$N = 365$, $n = 23$: $\frac{n(n-1)}{2N} = \frac{23 \cdot 22}{730} \approx 0{,}6932$,
$1 - e^{-0{,}6932} \approx 0{,}5000$, а точное значение
$1 - \prod_{i=1}^{22}(1 - i/365) \approx 0{,}5073$. Уже в группе из $23$ человек
с вероятностью больше $1/2$ у двоих совпадают дни рождения, хотя $23 \ll 365$.
\end{example}

\begin{theorem}[{о коллизиях, два списка; \cite[разд.~5.4]{hps}}]\label{l3:collision}
Пусть $A \subseteq S$, $\abs{A} = n$ (первый список из $n$ различных элементов), а
$b_1, \dots, b_m \in S$ выбраны независимо и равновероятно (второй список). Тогда
\[
  \Pr(\exists\, j:\ b_j \in A) = 1 - \Bigl(1 - \frac nN\Bigr)^m \ \ge\ 1 - e^{-nm/N}.
\]
\end{theorem}

\begin{proof}
Каждый $b_j$ не попадает в $A$ с вероятностью $1 - n/N$; события независимы, поэтому
ни один не попадает с вероятностью $(1 - n/N)^m \le (e^{-n/N})^m = e^{-nm/N}$
по лемме~\ref{l3:exp}. «Модель урны»: в урне $N$ шаров, из них $n$ красных; вынимаем
шар $m$ раз с возвращением "--- вероятность ни разу не вынуть красный равна
$(1 - n/N)^m$.
\end{proof}

\begin{corollary}[длина списков]\label{l3:lengths}
Пусть $0 < P < 1$ и $L = \ln\frac{1}{1 - P}$.
\begin{enumerate}
  \item Два списка: если $nm \ge N L$ (в частности, $n = m \ge \sqrt{N L}$), то коллизия
    между списками есть с вероятностью $\ge P$.
  \item Один список: если $n(n - 1) \ge 2 N L$, т.\,е.
    $n \ge \frac{1 + \sqrt{1 + 8NL}}{2} \approx \sqrt{2NL}$, то коллизия есть
    с вероятностью $\ge P$.
\end{enumerate}
При $P = \frac12$ ($L = \ln 2$): $n = m \approx 0{,}83\sqrt N$ для двух списков и
$n \approx 1{,}18\sqrt N$ для одного. При $P = 1 - e^{-9} \approx 0{,}9999$ хватает
$n = m = 3\sqrt N$.
\end{corollary}

\begin{proof}
$1 - e^{-nm/N} \ge P \iff e^{-nm/N} \le 1 - P \iff nm/N \ge L$; аналогично
$1 - e^{-n(n-1)/2N} \ge P \iff n^2 - n - 2NL \ge 0$, и остаётся решить квадратное неравенство.
\end{proof}

Итак, число элементов, которое нужно перебрать до коллизии, имеет порядок $\sqrt N$, а не $N$.
Это объясняет, например, почему у хеш-функции с $k$-битным значением стойкость к
поиску коллизий "--- лишь $2^{k/2}$ операций.

\taskexample{collision-birthday}

\begin{theorem}[ожидаемое время до первой коллизии]\label{l3:expected}
Пусть элементы $a_1, a_2, \dots \in S$ выбираются независимо и равновероятно, и $X$ "--- номер
первого элемента, совпавшего с одним из предыдущих. Тогда
$\mathsf{E}X \le \sqrt{\pi N/2} + 2$.
\end{theorem}

\begin{proof}
$X > k$ означает, что $a_1, \dots, a_k$ различны, поэтому по теореме~\ref{l3:birthday}
$\Pr(X > k) = \prod_{i=1}^{k-1}(1 - i/N) \le e^{-k(k-1)/2N} \le e^{-(k-1)^2/2N}$ при $k \ge 1$.
Для случайной величины с натуральными значениями $\mathsf{E}X = \sum_{k \ge 0} \Pr(X > k)$,
и
\[
  \mathsf{E}X \le 1 + \sum_{j \ge 0} e^{-j^2/2N}
  \le 2 + \int_0^\infty e^{-t^2/2N}\,dt = 2 + \sqrt{\frac{\pi N}{2}},
\]
где мы оценили сумму убывающей функции по $j \ge 1$ интегралом.
\end{proof}

Более точно, $\mathsf{E}X \sim \sqrt{\pi N/2} \approx 1{,}25\sqrt N$ \cite[разд.~14.1]{galbraith}.

\checkyourself{Сколько случайных $64$-битных значений хеш-функции нужно вычислить, чтобы
с вероятностью $1/2$ найти коллизию? Почему для DLP в группе порядка $N \approx 2^{256}$
говорят о «$128$-битной стойкости»?}

\section{Коллизионный алгоритм для DLP}

Пусть $\ord g = N$ и $h \in \langle g \rangle$. Если $g^{y} = h g^{z}$, то
$h = g^{y - z}$, т.\,е. $x \equiv y - z \pmod N$. Будем искать такое совпадение среди
случайных степеней~\cite[разд.~5.4]{hps}.

\begin{algorithm}[коллизионный алгоритм для DLP]
\begin{algorithmic}[1]
\Require $g$, $\ord g = N$, $h \in \langle g \rangle$; длина списков $n$.
\Ensure $x$: $g^x = h$.
\State выбрать различные случайные $y_1, \dots, y_n \in \Zn{N}$, вычислить $g^{y_i}$
  и сохранить пары $(g^{y_i}, y_i)$ в хеш-таблице
\State выбрать случайное $z \in \Zn{N}$, $w \gets h \cdot g^{z}$
\While{$w \notin \{g^{y_1}, \dots, g^{y_n}\}$}
  \State выбрать новое случайное $z \in \Zn{N}$, $w \gets h \cdot g^{z}$
\EndWhile
\State найти $i$ с $w = g^{y_i}$ и \Return $x = (y_i - z) \bmod N$
\end{algorithmic}
\end{algorithm}

Корректность очевидна. Элементы $g^{y_i}$ различны (различны показатели $y_i \bmod N$), а
каждое $h g^{z} = g^{x + z}$ равномерно распределено в $\langle g \rangle$ и не зависит от
предыдущих. Значит, к первому списку и к элементам $h g^{z}$ применима теорема~\ref{l3:collision}:
при $n$ элементах первого списка и $m$ шагах цикла коллизии ещё нет с вероятностью
$(1 - n/N)^m \le e^{-nm/N}$.

\begin{theorem}\label{l3:collision-dlp}
При $n = \ceil{\sqrt N}$ ожидаемое число шагов цикла не превосходит $\sqrt N$, и алгоритм
выполняет в среднем $\bigO(\sqrt N)$ возведений в степень, т.\,е. $\bigO(\sqrt N \log N)$
умножений в группе, и хранит $\bigO(\sqrt N)$ элементов. С вероятностью $\ge 1 - e^{-9}$
алгоритм завершится за $9\sqrt N$ шагов цикла.
\end{theorem}

\begin{proof}
Каждый шаг цикла успешен независимо от остальных с вероятностью $\pi = n/N \ge 1/\sqrt N$,
т.\,е. число шагов $m$ имеет геометрическое распределение и
$\mathsf{E} m = 1/\pi \le \sqrt N$. Первый этап "--- $n \approx \sqrt N$ возведений
в степень; на каждом шаге цикла "--- одно возведение в степень $\bigO(\log N)$ умножений
(лекция~6 первого семестра) и поиск в хеш-таблице за $\bigO(1)$ в среднем. Наконец, по
теореме~\ref{l3:collision} $\Pr(m > 9\sqrt N) \le e^{-n \cdot 9\sqrt N/N} \le e^{-9}$.
\end{proof}

Мы получили то же, что и в алгоритме Шенкса, но с вероятностным результатом и той же
памятью $\sqrt N$. Выигрыш в другом: идея «ищем совпадение двух случайных элементов» допускает
реализацию \emph{без памяти}, к которой мы и переходим.

\taskexample{collision-lists}

\checkyourself{Почему в первом списке показатели $y_i$ нужно брать различными, а во втором
списке "--- не обязательно? Как изменится оценка, если первый список содержит $n$ элементов,
а памяти хватает лишь на $n = \sqrt N / 100$ элементов?}

\section{Орбиты отображения конечного множества}

Пусть $S$ "--- конечное множество, $\abs{S} = N$, $f\colon S \to S$ "--- произвольное
отображение и $x_0 \in S$. Положим $x_{i+1} = f(x_i)$, т.\,е. $x_i = f^{i}(x_0)$.

\begin{definition}
\emph{Прямой орбитой} элемента $x_0$ называется множество
$\mathcal{O}^+(x_0) = \{x_0, f(x_0), f^2(x_0), \dots\} = \{x_i : i \ge 0\}$.
\end{definition}

Поскольку $S$ конечно, среди $x_0, x_1, \dots, x_N$ есть совпадающие. Пусть $x_j$ "--- первый
элемент последовательности, совпавший с одним из предыдущих: $x_j = x_T$, $T < j$; положим
$M = j - T$. Число $T \ge 0$ называется \emph{предпериодом} (длиной «хвоста»), $M \ge 1$
"--- \emph{периодом} (длиной цикла) орбиты.

\begin{theorem}[структура орбиты]\label{l3:orbit}
\begin{enumerate}
  \item Элементы $x_0, x_1, \dots, x_{T+M-1}$ попарно различны, и
    $\mathcal{O}^+(x_0) = \{x_0, \dots, x_{T+M-1}\}$; в частности, $T + M \le N$.
  \item $x_{i+M} = x_i$ при всех $i \ge T$.
  \item $x_i = x_j$ тогда и только тогда, когда $i = j$ или $i, j \ge T$ и
    $i \equiv j \pmod M$.
\end{enumerate}
\end{theorem}

\begin{proof}
1. Различие $x_0, \dots, x_{T+M-1}$ "--- это определение $j = T + M$ как первого повтора.

2. Индукция по $i \ge T$: $x_{T+M} = x_T$ по определению, и если $x_{i+M} = x_i$, то
$x_{i+1+M} = f(x_{i+M}) = f(x_i) = x_{i+1}$. Отсюда $x_{i + kM} = x_i$ при всех $k \ge 0$,
$i \ge T$, и каждый $x_i$ с $i \ge T + M$ совпадает с одним из $x_T, \dots, x_{T+M-1}$:
это доказывает и равенство в п.~1.

3. «$\Leftarrow$» следует из п.~2. «$\Rightarrow$»: пусть $x_i = x_j$, $i < j$. Каждый элемент
с индексом $k \ge T$ совпадает с $x_{r(k)}$, где $r(k) = T + ((k - T) \bmod M)$, а
$r(k) = k$ при $k < T$. Значит, $x_{r(i)} = x_{r(j)}$, где $r(i), r(j) \le T + M - 1$; по п.~1
$r(i) = r(j)$. Если $i < T$, то $r(i) = i < T$, а $r(j) \ge T$ при $j \ge T$ и $r(j) = j \ne i$
при $j < T$ "--- противоречие. Значит, $i, j \ge T$ и $(i - T) \bmod M = (j - T) \bmod M$,
т.\,е. $i \equiv j \pmod M$.
\end{proof}

Отсюда название: если нарисовать орбиту, «хвост» $x_0, \dots, x_{T-1}$ переходит в цикл
$x_T, \dots, x_{T+M-1}$, и картинка похожа на греческую букву $\rho$
(рис.~\ref{l3:fig-rho}).

\begin{figure}[ht]
\centering
\rhoshape[scale=0.8]{3}{4}
\caption{Орбита $x_0 = 0$ при $f(x) = x^2 + 2 \bmod 17$: $0, 2, 6, \mid 4, 1, 3, 11, \mid 4, \dots$;
предпериод $T = 3$, период $M = 4$.}\label{l3:fig-rho}
\end{figure}

\begin{example}
$S = \Zn{17}$, $f(x) = x^2 + 2$, $x_0 = 0$: $x_1 = 2$, $x_2 = 6$, $x_3 = 38 \bmod 17 = 4$,
$x_4 = 18 \bmod 17 = 1$, $x_5 = 3$, $x_6 = 11$, $x_7 = 123 \bmod 17 = 4 = x_3$. Значит,
$T = 3$, $M = 4$, и, например, $x_{10} = x_6$, так как $10 \equiv 6 \pmod 4$ и $6, 10 \ge 3$.
\end{example}

\subsection{Поиск цикла: алгоритм Флойда}

Чтобы найти совпадение $x_i = x_j$, можно хранить все $x_i$ "--- но это снова $\bigO(\sqrt N)$
памяти. Р.~Флойд заметил, что достаточно сравнивать $x_i$ с $x_{2i}$.

\begin{theorem}[Флойд]\label{l3:floyd}
Существует $i$, $1 \le i \le T + M$, для которого $x_i = x_{2i}$. Наименьшее такое
$i$ "--- это наименьшее кратное $M$, не меньшее $\max(T, 1)$.
\end{theorem}

\begin{proof}
При $i \ge 1$ имеем $2i \ne i$, и по теореме~\ref{l3:orbit} $x_i = x_{2i}$ $\iff$
$i \ge T$ и $2i \equiv i \pmod M$, т.\,е. $M \mid i$. Среди $M$ подряд идущих чисел
$\max(T,1), \dots, \max(T,1) + M - 1$ ровно одно делится на $M$; оно не превосходит
$\max(T, 1) + M - 1 \le T + M$.
\end{proof}

\begin{algorithm}[Флойд, «черепаха и заяц»]
\begin{algorithmic}[1]
\Require $f\colon S \to S$, $x_0 \in S$.
\Ensure $i \ge 1$ с $x_i = x_{2i}$.
\State $u \gets f(x_0)$, \quad $v \gets f(f(x_0))$, \quad $i \gets 1$
  \Comment{$u = x_i$, $v = x_{2i}$}
\While{$u \ne v$}
  \State $u \gets f(u)$, \quad $v \gets f(f(v))$, \quad $i \gets i + 1$
\EndWhile
\State \Return $i$
\end{algorithmic}
\end{algorithm}

Алгоритм хранит два элемента $S$ и вычисляет $f$ ровно $3i \le 3(T + M)$ раз.

\checkyourself{Найдите для орбиты из примера ($f(x) = x^2 + 2 \bmod 17$, $x_0 = 0$) наименьшее
$i$ с $x_i = x_{2i}$ и проверьте теорему Флойда. Может ли случиться, что $i = T + M$?}

\subsection{Длина хвоста и цикла случайного отображения}

Насколько длинна орбита? Для «плохого» $f$ (например, циклической перестановки $S$)
$T + M = N$. Но если $f$ выбрано \emph{случайно} среди всех $N^N$ отображений $S \to S$,
то значения $f$ в разных точках независимы и равномерно распределены. Пока точки
$x_0, \dots, x_k$ различны, значения $x_{k+1} = f(x_k)$ ещё «не открыты», поэтому
$x_1, x_2, \dots$ до первого повтора ведут себя как независимые случайные элементы $S$.
Номер первого повтора равен $T + M$, и по теореме~\ref{l3:expected} (с учётом $x_0$)
$\mathsf{E}(T + M) \le \sqrt{\pi N/2} + 2$.

\begin{theorem}[{Харрис; \emph{без доказательства}, см.~\cite[разд.~14.1]{galbraith}}]\label{l3:random-map}
Для случайного отображения $f\colon S \to S$ и случайного $x_0$
\[
  \mathsf{E}\,T \sim \mathsf{E}\,M \sim \sqrt{\frac{\pi N}{8}} \approx 0{,}627\sqrt N, \qquad
  \mathsf{E}(T + M) \sim \sqrt{\frac{\pi N}{2}} \approx 1{,}253\sqrt N.
\]
Ожидаемое число итераций алгоритма Флойда равно $\frac{\pi^2}{12}\sqrt{\frac{\pi N}{2}}
\approx 1{,}03\sqrt N$.
\end{theorem}

Эвристически деление поровну понятно: первый повтор $x_{T+M}$ попадает в «равномерно
случайную» точку уже пройденного пути $x_0, \dots, x_{T+M-1}$, поэтому в среднем $T \approx M$.

В приложениях $f$, разумеется, не случайно, а задано формулой. Основная \emph{эвристика}
метода Полларда: «хорошо перемешивающие» отображения (например, $x \mapsto x^2 + 1 \bmod n$
или итерация из следующего раздела) ведут себя как случайные, и оценки
теоремы~\ref{l3:random-map} для них выполняются. Это подтверждается экспериментами, хотя
и не доказано.

\begin{remark}[вариант Брента]
Р.~Брент предложил хранить $y = x_{2^k - 1}$ и сравнивать с ним $x_j$ при
$2^k \le j < 2^{k+1}$, затем заменять $y$ на $x_{2^{k+1}-1}$. Совпадение найдётся,
когда $2^k - 1 \ge T$ и $2^k \ge M$, т.\,е. не позднее, чем через $\bigO(T + M)$ шагов,
причём на каждом шаге $f$ вычисляется \emph{один} раз (у Флойда "--- три). На практике
это заметно (примерно на треть) уменьшает число вычислений $f$ \cite[гл.~14]{galbraith}.
\end{remark}

\checkyourself{Почему в варианте Брента достаточно условий $2^k - 1 \ge T$ и $2^k \ge M$?
Какое из них нарушается, если хвост очень длинный?}

\section{\texorpdfstring{$\rho$}{ρ}-алгоритм Полларда для DLP}

Пусть $G = \langle g \rangle \le \Znx{p}$, $\ord g = n$ (обычно $n = q$ "--- простое число),
$h \in G$, и нужно найти $x$ с $g^x = h$. Будем строить орбиту отображения $f\colon G \to G$,
для элементов которой \emph{известно представление} $x_i = g^{\alpha_i} h^{\beta_i}$.

Разобьём $G$ на три примерно равные части по остатку от деления (представителя
$x \in \{1, \dots, p - 1\}$) на $3$: $S_k = \{x \in G : x \equiv k \pmod 3\}$, и положим
\[
  f(x) = \begin{cases}
    x^2, & x \in S_0,\\
    h x, & x \in S_1,\\
    g x, & x \in S_2,
  \end{cases}
  \qquad
  (\alpha, \beta) \mapsto \begin{cases}
    (2\alpha, 2\beta), & x \in S_0,\\
    (\alpha, \beta + 1), & x \in S_1,\\
    (\alpha + 1, \beta), & x \in S_2,
  \end{cases}
\]
(показатели "--- по модулю $n$), начиная с $x_0 = g^{\alpha_0} h^{\beta_0}$ для случайных
$\alpha_0, \beta_0$ (или $x_0 = 1$, $\alpha_0 = \beta_0 = 0$).

\begin{lemma}\label{l3:invariant}
При всех $i \ge 0$ выполняется $x_i = g^{\alpha_i} h^{\beta_i}$.
\end{lemma}

\begin{proof}
Индукция: при $i = 0$ "--- по выбору $x_0$. Если $x_i = g^{\alpha_i} h^{\beta_i}$, то
$x_i^2 = g^{2\alpha_i} h^{2\beta_i}$, $h x_i = g^{\alpha_i} h^{\beta_i + 1}$,
$g x_i = g^{\alpha_i + 1} h^{\beta_i}$, а показатели можно приводить по модулю $n$,
так как $g^n = h^n = 1$.
\end{proof}

Умножения на $g$ и $h$ «перемешивают» орбиту, а возведение в квадрат нужно, чтобы показатели
менялись нелинейно: при одних лишь умножениях на $g$ и $h$ отображение было бы «слишком
регулярным» и не похожим на случайное.

Найдя алгоритмом Флойда $i$ с $x_i = x_{2i}$, получаем
$g^{\alpha_i} h^{\beta_i} = g^{\alpha_{2i}} h^{\beta_{2i}}$, т.\,е. при $h = g^x$
\begin{equation}\label{l3:congruence}
  (\beta_i - \beta_{2i})\, x \equiv \alpha_{2i} - \alpha_i \pmod n.
\end{equation}

\begin{lemma}[линейное сравнение]\label{l3:linear}
Пусть $u, v \in \ZZ$, $d = \gcd(u, n)$. Сравнение $u x \equiv v \pmod n$ разрешимо тогда
и только тогда, когда $d \mid v$, и в этом случае имеет ровно $d$ решений по модулю $n$:
\[
  x \equiv \tilde x + k\,\frac nd \pmod n, \quad k = 0, 1, \dots, d - 1, \qquad
  \tilde x \equiv \frac vd \Bigl(\frac ud\Bigr)^{-1} \pmod{\frac nd}.
\]
\end{lemma}

\begin{proof}
Если $ux - v = tn$, то $d \mid ux - tn = v$. Обратно, пусть $d \mid v$. Тогда
$ux \equiv v \pmod n$ $\iff$ $\frac ud x \equiv \frac vd \pmod{\frac nd}$ (делим равенство
$ux - v = tn$ на $d$), а $\gcd(\frac ud, \frac nd) = 1$, поэтому $\frac ud$ обратим по модулю
$\frac nd$ (расширенный алгоритм Евклида), и решение единственно по модулю $\frac nd$:
$x \equiv \tilde x \pmod{\frac nd}$. По модулю $n$ этому классу соответствуют ровно
$d$ различных вычетов $\tilde x + k \frac nd$, $0 \le k < d$.
\end{proof}

Применяя лемму к~\eqref{l3:congruence} с $u = \beta_i - \beta_{2i}$, $v = \alpha_{2i} - \alpha_i$:
\begin{itemize}
  \item если $n = q$ простое и $u \not\equiv 0$, то $d = 1$ и $x$ находится однозначно;
  \item если $d > 1$ (возможно лишь при составном $n$), получаем $d$ кандидатов
    $\tilde x + k\frac nd$, и верный находится проверкой $g^{x} = h$ ($d$ возведений
    в степень; при больших $d$ "--- с помощью DLP в подгруппе порядка $d$ для
    $g^{n/d}$, лекция~2);
  \item если $u \equiv 0 \pmod n$ (и тогда $v \equiv 0$), совпадение бесполезно "--- начинаем
    заново с другими $\alpha_0, \beta_0$. При простом $n = q$ это случается с вероятностью
    $\approx 1/q$.
\end{itemize}
Сравнение~\eqref{l3:congruence} разрешимо всегда: $x = \log_g h$ ему удовлетворяет.

\begin{algorithm}[$\rho$-алгоритм Полларда для DLP]
\begin{algorithmic}[1]
\Require $g$, $\ord g = n$, $h \in \langle g \rangle$.
\Ensure $x$: $g^x = h$.
\State выбрать случайные $\alpha_0, \beta_0 \in \Zn{n}$, $x_0 \gets g^{\alpha_0}h^{\beta_0}$
\State $(u, a, b) \gets (x_0, \alpha_0, \beta_0)$, \quad $(v, A, B) \gets (x_0, \alpha_0, \beta_0)$
\State $(u, a, b) \gets F(u, a, b)$, \quad $(v, A, B) \gets F(F(v, A, B))$
  \Comment{$F$ "--- шаг $f$ с пересчётом показателей}
\While{$u \ne v$}
  \State $(u, a, b) \gets F(u, a, b)$, \quad $(v, A, B) \gets F(F(v, A, B))$
\EndWhile
\If{$b \equiv B \pmod n$} \State начать заново (строка 1) \EndIf
\State $d \gets \gcd(b - B, n)$; решить $(b - B)x \equiv A - a \pmod n$ (лемма~\ref{l3:linear})
\State \Return тот из $d$ кандидатов, для которого $g^x = h$
\end{algorithmic}
\end{algorithm}

\begin{theorem}\label{l3:rho-complexity}
В предположении, что $f$ ведёт себя как случайное отображение группы $G$ порядка $n$,
$\rho$-алгоритм Полларда находит коллизию за ожидаемое число $\approx 3 \cdot 1{,}03\sqrt n$
групповых операций (плюс $\bigO(\log^2 n)$ на решение сравнения) и хранит $\bigO(1)$
элементов группы и показателей.
\end{theorem}

\begin{proof}
По теореме~\ref{l3:floyd} алгоритм Флойда завершается при $i \le T + M$, вычисляя $f$ $3i$
раз, а по теореме~\ref{l3:random-map} $\mathsf{E}\, i \approx 1{,}03\sqrt n$ (а грубая
оценка $\mathsf{E}\, i \le \mathsf{E}(T + M) \le \sqrt{\pi n/2} + 2$ доказана выше).
Каждое вычисление $f$ "--- одно умножение в $G$ и одно-два сложения/удвоения по модулю $n$.
Хранятся только $(u, a, b)$ и $(v, A, B)$. Решение сравнения "--- расширенный алгоритм Евклида
и $d$ проверок.
\end{proof}

Вместе с алгоритмом Полига"--~Хеллмана получаем: DLP в группе порядка
$n = \prod q_j^{e_j}$ решается за $\bigO\bigl(\sum_j e_j(\log n + \sqrt{q_j})\bigr)$ групповых
операций с \emph{константной} памятью. Поэтому порядок группы в криптографии должен иметь
простой делитель $q \gtrsim 2^{2k}$ для $k$-битной стойкости.

\begin{remark}
Отображение с тремя частями "--- исторически первое; Э.~Теске показала, что оно ведёт себя
несколько хуже случайного, и предложила «$r$-аддитивные» блуждания
$x \mapsto x \cdot g^{a_j} h^{b_j}$ при $x \in S_j$, $j = 1, \dots, r$ ($r \approx 20$),
которые практически неотличимы от случайных \cite[разд.~14.2]{galbraith}.
\end{remark}

\taskexample{rho-floyd}

В следующем примере порядок $n = 36$ основания составной, и сравнение имеет
$d = 3$ решения, из которых нужное выбирается проверкой.

\taskexample{rho-composite}

\checkyourself{Почему при простом $n = q$ всегда $d = 1$ (если $\beta_i \not\equiv \beta_{2i}$)?
Сколько кандидатов может получиться, если $g$ "--- первообразный корень по модулю $p$?}

\checkyourself{Что произойдёт, если в отображении $f$ убрать ветвь $x \mapsto x^2$? А если
начать с $x_0 = 1$ и $1 \in S_0$?}

\section{Метод выделенных точек и параллельный поиск}

Алгоритм Флойда плохо распараллеливается: $m$ независимых запусков на $m$ вычислителях
ускоряют поиск лишь примерно в $\sqrt m$ раз (каждый запуск ищет коллизию \emph{внутри своей}
орбиты, а минимум из $m$ времён порядка $\sqrt n$ имеет порядок $\sqrt{n/m}$). Нужен способ
находить коллизии \emph{между} орбитами разных вычислителей, не храня все точки.

\begin{definition}
Зафиксируем легко проверяемое свойство элементов $G$, которым обладает доля $\theta$ всех
элементов, например: «$k$ младших двоичных разрядов представления $x$ нулевые» (тогда
$\theta \approx 2^{-k}$). Элементы с этим свойством называются \emph{выделенными точками}.
\end{definition}

\begin{algorithm}[$\rho$-метод с выделенными точками, ван~Ооршот"--~Винер]
\begin{algorithmic}[1]
\Require $g$, $\ord g = n$, $h$; $m$ вычислителей, свойство «выделенности».
\Ensure $x$: $g^x = h$.
\State \textbf{Каждый вычислитель:} выбрать случайные $\alpha, \beta$,
  $y \gets g^{\alpha}h^{\beta}$
\While{сервер не дал сигнал остановки}
  \State $(y, \alpha, \beta) \gets F(y, \alpha, \beta)$
  \If{$y$ "--- выделенная точка}
    \State отправить $(y, \alpha, \beta)$ на сервер
    \State (иногда) начать новое блуждание со случайной точки
  \EndIf
\EndWhile
\State \textbf{Сервер:} хранит выделенные точки; получив $(y, \alpha', \beta')$ при уже
  сохранённой $(y, \alpha, \beta)$ с $\beta \not\equiv \beta'$, решает
  $(\beta - \beta')x \equiv \alpha' - \alpha \pmod n$ и останавливает всех.
\end{algorithmic}
\end{algorithm}

Почему это работает: все вычислители итерируют \emph{одно и то же} отображение $f$. Как только
два блуждания (или одно блуждание само с собой) попадут в одну точку, дальше они идут
\emph{одинаково} и, следовательно, придут в одну и ту же выделенную точку "--- там коллизию
и обнаружит сервер. Если выделенных точек на цикле нет, блуждание зацикливается навсегда;
поэтому блуждание, сделавшее, скажем, $20/\theta$ шагов без выделенных точек, прерывают и
начинают заново.

\begin{theorem}\label{l3:dp}
Пусть $f$ ведёт себя как случайное отображение, а каждая новая точка является выделенной
независимо с вероятностью $\theta$. Тогда при $m$ вычислителях ожидаемое время (число шагов
каждого вычислителя) до обнаружения коллизии примерно
\[
  \frac{\sqrt{\pi n/2}}{m} + \frac{1}{\theta},
\]
а сервер хранит в среднем около $\theta\sqrt{\pi n/2}$ точек.
\end{theorem}

\begin{proof}[Обоснование (эвристическое)]
Пока коллизии нет, все вычисленные точки всех блужданий различны и ведут себя как
независимые случайные элементы $G$ (новые значения случайного $f$). Поэтому по
теореме~\ref{l3:expected} первая коллизия наступает в среднем после
$\approx \sqrt{\pi n/2}$ вычисленных \emph{всеми} вычислителями точек, т.\,е. через
$\approx \sqrt{\pi n/2}/m$ шагов каждого. После коллизии блуждание идёт по уже пройденному
пути до первой выделенной точки; число шагов до неё имеет геометрическое распределение
с параметром $\theta$, и его математическое ожидание
$\sum_{k \ge 1} k\theta(1 - \theta)^{k-1} = 1/\theta$. Среди $\approx \sqrt{\pi n/2}$
вычисленных точек выделенных в среднем $\theta\sqrt{\pi n/2}$ "--- их и хранит сервер.
\end{proof}

Таким образом, ускорение \emph{линейно} по $m$, а параметр $\theta$ позволяет выбирать
компромисс: малое $\theta$ "--- мало памяти, но долго идти до выделенной точки после
коллизии; при $1/\theta \ll \sqrt n/m$ вторым слагаемым можно пренебречь. Например, при
$n \approx 2^{100}$, $m = 2^{20}$ и $\theta = 2^{-20}$ каждый вычислитель делает $\approx 2^{30}$
шагов, а сервер хранит $\approx 2^{30}$ точек. Именно так решены рекордные задачи ECDLP
(например, для кривых над полями из $\approx 2^{112}$ элементов).

\taskexample{rho-distinguished}

\checkyourself{Почему выделенные точки должны определяться свойством самого элемента группы,
а не номером шага блуждания? Что произойдёт при $\theta = 1$ и при $\theta \to 0$?}

\section{Другие применения}

\begin{itemize}
  \item \textbf{$\rho$-алгоритм Полларда для факторизации}~\cite{pollard-rho-factor}. Пусть
    $N$ составное, $p \mid N$ "--- неизвестный простой делитель, $f(x) = x^2 + 1 \bmod N$.
    Последовательность $x_i \bmod p$ "--- орбита отображения $\Zn{p} \to \Zn{p}$, и через
    $\bigO(\sqrt p)$ шагов $x_i \equiv x_{2i} \pmod p$, хотя, как правило,
    $x_i \ne x_{2i}$. Коллизию по неизвестному модулю $p$ обнаруживает
    $\gcd(x_{2i} - x_i, N) > 1$. Время "--- $\bigO(\sqrt p) = \bigO(N^{1/4})$ умножений по
    модулю $N$ (о факторизации "--- в лекциях~5 и~6).
  \item \textbf{ECDLP}: $\rho$-метод "--- лучший известный алгоритм для DLP на эллиптических кривых
    общего вида (лекция~9); поэтому для $128$-битной стойкости берут кривые порядка
    $\approx 2^{256}$.
  \item \textbf{Кенгуру ($\lambda$-метод) Полларда}: если известно, что $x$ лежит в отрезке длины $w$,
    тот же приём с выделенными точками даёт время $\bigO(\sqrt w)$
    \cite[разд.~14.5]{galbraith}.
  \item \textbf{Коллизии хеш-функций и атака «встреча посередине»} используют те же оценки
    теоремы о коллизиях.
\end{itemize}

\checkyourselfsection

\begin{furtherreading}
  \readitem[разд.~5.4--5.5]{hps}{парадокс дней рождения, теорема о коллизиях, коллизионный
    алгоритм для DLP, $\rho$-алгоритм Полларда и алгоритм Флойда (2-е изд.)}
  \readitem[гл.~14, разд.~14.1--14.4]{galbraith}{случайные отображения, $\rho$-алгоритм, циклы
    Флойда и Брента, $r$-аддитивные блуждания, выделенные точки и параллельный поиск}
  \readitem[разд.~3.6.3]{hac}{$\rho$-алгоритм Полларда для DLP (с той же разбивкой на три
    части)}
  \readitem[разд.~5.2.2]{crandall-pomerance}{$\rho$-метод для дискретного логарифма}
  \readitem{pollard-rho-dlp}{оригинальная статья Полларда}
  \readitem{van-oorschot-wiener}{параллельный поиск коллизий с выделенными точками}
  \readitem[лекция «Generic algorithms for the discrete logarithm problem»]{mit-18783}{конспекты: универсальные алгоритмы DLP, $\rho$-метод
    и его анализ}
\end{furtherreading}
\end{lecture}
